% T2 Loss terms and weights (Method 4.2). Owner: Method writer.
% Sources: defaults loss.py:436-452 (add_specific_args); terms loss.py:314-411 (forward);
% Final and Control: slurm/experiments/exp7_final_model.sh:210-222 (ARM fix / control), :759-771 (train.py command);
% Student: slurm/experiments/exp8_distill_vitb.sh header ("targets", "KD") and :319-331;
% distillation terms train.py:505-539; hand mask loss.py:231-253 (38 bones in the smoke log, H 0a step 6);
% shape mask loss.py:255-310 (off in all three columns; on in Final-masked, job 3612325).
% Boost renormalisation x4 -> x2.41 / x0.60: H 1 "x4 boost"; loss.py configure_joint_weights.
\begin{table}[!tbp]
\centering\footnotesize
\setlength{\tabcolsep}{3pt}
\begin{tabular}{@{}>{\raggedright}p{2.6cm}>{\raggedright}p{5.35cm}cccc@{}}
\toprule
Term (log name; \texttt{loss.py}) & What it compares, in plain words & Default & Final & Control & Student \tabularnewline
\midrule
Detection (\texttt{bce}; 314) & CornerNet focal loss on the per-token score ``an anchor joint lies in this token'' & \MethodWDet & \MethodWDet & \MethodWDet & \MethodWDet \tabularnewline
Offset (\texttt{offset}; 315) & L1 on the anchor's position inside its token & \MethodWOffset & \MethodWOffset & \MethodWOffset & \MethodWOffset \tabularnewline
Rotations (\texttt{rotmat}; 319--335) & L1 over the 9 entries of every bone's parent-relative rotation matrix & \MethodWRot & \MethodWRot & \MethodWRot & \MethodWRot \tabularnewline
\quad neck and hand boost & multiplier on neck, head and hand bones, then renormalised to mean 1 ($\times$\MethodBoostOn{} for them, $\times$\MethodBoostOff{} for the rest) & 1 & \MethodBoost & \MethodBoost & \MethodBoost \tabularnewline
Shape (\texttt{shape}; 336--343) & L1 on the \MethodNumPheno{} phenotypes after the sigmoid & \MethodWShapeDefault & \MethodWShape & \MethodWShape & \MethodWShape \tabularnewline
Distance (\texttt{dist}; 344) & L1 on the anchor's log-depth, normalised by the focal length & \MethodWDistDefault & \MethodWDist & \MethodWDistDefault & \MethodWDist \tabularnewline
Translation (\texttt{transl}; 345) & L1 on the anchor's 3D position in the camera frame & \MethodWDistDefault & \MethodWDist & \MethodWDistDefault & \MethodWDist \tabularnewline
3D joints (\texttt{j3d}; 367--377) & L1 on the \MethodNumBones{} joints, each skeleton moved to its own pelvis & \MethodWThreeD & \MethodWThreeD & \MethodWThreeD & \MethodWThreeD \tabularnewline
3D vertices (\texttt{v3d}; 378) & the same for the mesh vertices & \MethodWThreeD & \MethodWThreeD & \MethodWThreeD & \MethodWThreeD \tabularnewline
2D joints (\texttt{j2d}; 383--384) & L1 in pixels on the projected joints, placed by the predicted anchor and depth & \MethodWTwoDDefault & \MethodWTwoD & \MethodWTwoD & \MethodWTwoD \tabularnewline
2D vertices (\texttt{v2d}; 381--382) & the same for the projected vertices & \MethodWTwoDDefault & \MethodWTwoD & \MethodWTwoD & \MethodWTwoD \tabularnewline
\quad first epoch of the 2D terms & \texttt{-{}-start\_2d\_epoch} & \MethodStartTwoDDefault & \MethodStartTwoD & \MethodStartTwoD & \MethodStartTwoD \tabularnewline
Global rotation (\texttt{global\_rot}; 350--363) & L1 on the camera-frame orientation of head, neck03 and both wrists & 0 & 0 & 0 & 0 \tabularnewline
Feature distillation (\texttt{kd}) & KL divergence between the softened channel distributions of student and teacher tokens, $T=\MethodKDTemp$ & -- & -- & -- & \MethodKDLambda \tabularnewline
Output distillation (\texttt{kd\_out}) & L1 to the teacher's rotations, phenotypes and depth on the same people & -- & -- & -- & \MethodKDOutLambda \tabularnewline
\midrule
Joint the targets describe & \texttt{-{}-gt\_center\_joint}: transl, dist, detection token and offset targets & pelvis & head & pelvis & head \tabularnewline
Bone correction & constant output rotation on neck02, head, metacarpal1.L, finger4-1.R & none & on & none & on \tabularnewline
\bottomrule
\end{tabular}
\caption[Loss terms and their weights]{Loss terms and their weights. ``Default'' is the value in \texttt{loss.py}; Final, Control and Student are the
values of those runs (model names in Table~\ref{tab:model-names}). The total loss is the weighted sum of all rows above
the rule. For people from 3DPW the rotation and global-rotation terms skip the \MethodHandMaskBones{} finger, thumb
and metacarpal bones in all three runs; the mesh terms are not masked. The shape mask for 3DPW people is off in all
three runs and on only in Final-masked (Section~\ref{sec:method-real}). The Student keeps the teacher's bone correction
during training and is evaluated with its own, re-estimated after training (Section~\ref{sec:method-student}).}
\label{tab:loss-weights}
\end{table}
